% Options for packages loaded elsewhere
\PassOptionsToPackage{unicode}{hyperref}
\PassOptionsToPackage{hyphens}{url}
\documentclass[
]{article}
\usepackage{xcolor}
\usepackage[margin=1in]{geometry}
\usepackage{amsmath,amssymb}
\setcounter{secnumdepth}{-\maxdimen} % remove section numbering
\usepackage{iftex}
\ifPDFTeX
  \usepackage[T1]{fontenc}
  \usepackage[utf8]{inputenc}
  \usepackage{textcomp} % provide euro and other symbols
\else % if luatex or xetex
  \usepackage{unicode-math} % this also loads fontspec
  \defaultfontfeatures{Scale=MatchLowercase}
  \defaultfontfeatures[\rmfamily]{Ligatures=TeX,Scale=1}
\fi
\usepackage{lmodern}
\ifPDFTeX\else
  % xetex/luatex font selection
\fi
% Use upquote if available, for straight quotes in verbatim environments
\IfFileExists{upquote.sty}{\usepackage{upquote}}{}
\IfFileExists{microtype.sty}{% use microtype if available
  \usepackage[]{microtype}
  \UseMicrotypeSet[protrusion]{basicmath} % disable protrusion for tt fonts
}{}
\makeatletter
\@ifundefined{KOMAClassName}{% if non-KOMA class
  \IfFileExists{parskip.sty}{%
    \usepackage{parskip}
  }{% else
    \setlength{\parindent}{0pt}
    \setlength{\parskip}{6pt plus 2pt minus 1pt}}
}{% if KOMA class
  \KOMAoptions{parskip=half}}
\makeatother
\usepackage{color}
\usepackage{fancyvrb}
\newcommand{\VerbBar}{|}
\newcommand{\VERB}{\Verb[commandchars=\\\{\}]}
\DefineVerbatimEnvironment{Highlighting}{Verbatim}{commandchars=\\\{\}}
% Add ',fontsize=\small' for more characters per line
\usepackage{framed}
\definecolor{shadecolor}{RGB}{248,248,248}
\newenvironment{Shaded}{\begin{snugshade}}{\end{snugshade}}
\newcommand{\AlertTok}[1]{\textcolor[rgb]{0.94,0.16,0.16}{#1}}
\newcommand{\AnnotationTok}[1]{\textcolor[rgb]{0.56,0.35,0.01}{\textbf{\textit{#1}}}}
\newcommand{\AttributeTok}[1]{\textcolor[rgb]{0.13,0.29,0.53}{#1}}
\newcommand{\BaseNTok}[1]{\textcolor[rgb]{0.00,0.00,0.81}{#1}}
\newcommand{\BuiltInTok}[1]{#1}
\newcommand{\CharTok}[1]{\textcolor[rgb]{0.31,0.60,0.02}{#1}}
\newcommand{\CommentTok}[1]{\textcolor[rgb]{0.56,0.35,0.01}{\textit{#1}}}
\newcommand{\CommentVarTok}[1]{\textcolor[rgb]{0.56,0.35,0.01}{\textbf{\textit{#1}}}}
\newcommand{\ConstantTok}[1]{\textcolor[rgb]{0.56,0.35,0.01}{#1}}
\newcommand{\ControlFlowTok}[1]{\textcolor[rgb]{0.13,0.29,0.53}{\textbf{#1}}}
\newcommand{\DataTypeTok}[1]{\textcolor[rgb]{0.13,0.29,0.53}{#1}}
\newcommand{\DecValTok}[1]{\textcolor[rgb]{0.00,0.00,0.81}{#1}}
\newcommand{\DocumentationTok}[1]{\textcolor[rgb]{0.56,0.35,0.01}{\textbf{\textit{#1}}}}
\newcommand{\ErrorTok}[1]{\textcolor[rgb]{0.64,0.00,0.00}{\textbf{#1}}}
\newcommand{\ExtensionTok}[1]{#1}
\newcommand{\FloatTok}[1]{\textcolor[rgb]{0.00,0.00,0.81}{#1}}
\newcommand{\FunctionTok}[1]{\textcolor[rgb]{0.13,0.29,0.53}{\textbf{#1}}}
\newcommand{\ImportTok}[1]{#1}
\newcommand{\InformationTok}[1]{\textcolor[rgb]{0.56,0.35,0.01}{\textbf{\textit{#1}}}}
\newcommand{\KeywordTok}[1]{\textcolor[rgb]{0.13,0.29,0.53}{\textbf{#1}}}
\newcommand{\NormalTok}[1]{#1}
\newcommand{\OperatorTok}[1]{\textcolor[rgb]{0.81,0.36,0.00}{\textbf{#1}}}
\newcommand{\OtherTok}[1]{\textcolor[rgb]{0.56,0.35,0.01}{#1}}
\newcommand{\PreprocessorTok}[1]{\textcolor[rgb]{0.56,0.35,0.01}{\textit{#1}}}
\newcommand{\RegionMarkerTok}[1]{#1}
\newcommand{\SpecialCharTok}[1]{\textcolor[rgb]{0.81,0.36,0.00}{\textbf{#1}}}
\newcommand{\SpecialStringTok}[1]{\textcolor[rgb]{0.31,0.60,0.02}{#1}}
\newcommand{\StringTok}[1]{\textcolor[rgb]{0.31,0.60,0.02}{#1}}
\newcommand{\VariableTok}[1]{\textcolor[rgb]{0.00,0.00,0.00}{#1}}
\newcommand{\VerbatimStringTok}[1]{\textcolor[rgb]{0.31,0.60,0.02}{#1}}
\newcommand{\WarningTok}[1]{\textcolor[rgb]{0.56,0.35,0.01}{\textbf{\textit{#1}}}}
\usepackage{graphicx}
\makeatletter
\newsavebox\pandoc@box
\newcommand*\pandocbounded[1]{% scales image to fit in text height/width
  \sbox\pandoc@box{#1}%
  \Gscale@div\@tempa{\textheight}{\dimexpr\ht\pandoc@box+\dp\pandoc@box\relax}%
  \Gscale@div\@tempb{\linewidth}{\wd\pandoc@box}%
  \ifdim\@tempb\p@<\@tempa\p@\let\@tempa\@tempb\fi% select the smaller of both
  \ifdim\@tempa\p@<\p@\scalebox{\@tempa}{\usebox\pandoc@box}%
  \else\usebox{\pandoc@box}%
  \fi%
}
% Set default figure placement to htbp
\def\fps@figure{htbp}
\makeatother
\setlength{\emergencystretch}{3em} % prevent overfull lines
\providecommand{\tightlist}{%
  \setlength{\itemsep}{0pt}\setlength{\parskip}{0pt}}
\usepackage{bookmark}
\IfFileExists{xurl.sty}{\usepackage{xurl}}{} % add URL line breaks if available
\urlstyle{same}
\hypersetup{
  hidelinks,
  pdfcreator={LaTeX via pandoc}}

\author{}
\date{\vspace{-2.5em}}

\begin{document}

\section{Dữ liệu và thiết lập thực
nghiệm}\label{dux1eef-liux1ec7u-vuxe0-thiux1ebft-lux1eadp-thux1ef1c-nghiux1ec7m}

\subsection{Quy trình thực hiện}\label{quy-truxecnh-thux1ef1c-hiux1ec7n}

Quy trình thực nghiệm được thực hiện theo các bước sau:

• EDA dữ liệu trên toàn bộ dataset.

• Chia data thành train và test (validation set), thực hiện xây dựng mô
hình hồi quy tuyến tính và các phân tích trên tập train,

• Áp dụng các tiêu chuẩn lựa chọn biến gồm Adjusted \(R^2\), Mallows'
\(C_p\), AIC và BIC.

• Thực hiện các thủ tục lựa chọn biến gồm Forward Selection, Backward
Elimination và Stepwise Selection.

• So sánh các mô hình thu được từ các phương pháp lựa chọn biến.

• Đánh giá khả năng dự báo trên tập test thông qua RMSE, MAE và test
\(R^2\).

• Lựa chọn mô hình cuối cùng và đưa ra nhận xét.

\subsection{EDA}\label{eda}

\begin{verbatim}
## Loaded lars 1.3
\end{verbatim}

\begin{verbatim}
## corrplot 0.95 loaded
\end{verbatim}

\begin{Shaded}
\begin{Highlighting}[]
\FunctionTok{data}\NormalTok{(diabetes)}
\NormalTok{df }\OtherTok{\textless{}{-}} \FunctionTok{as.data.frame}\NormalTok{(}\FunctionTok{cbind}\NormalTok{(diabetes}\SpecialCharTok{$}\NormalTok{x, }\AttributeTok{target =}\NormalTok{ diabetes}\SpecialCharTok{$}\NormalTok{y))}
\FunctionTok{summary}\NormalTok{(df)}
\end{Highlighting}
\end{Shaded}

\begin{verbatim}
##       age                 sex                bmi                 map           
##  Min.   :-0.107226   Min.   :-0.04464   Min.   :-0.090275   Min.   :-0.112400  
##  1st Qu.:-0.037299   1st Qu.:-0.04464   1st Qu.:-0.034229   1st Qu.:-0.036656  
##  Median : 0.005383   Median :-0.04464   Median :-0.007284   Median :-0.005671  
##  Mean   : 0.000000   Mean   : 0.00000   Mean   : 0.000000   Mean   : 0.000000  
##  3rd Qu.: 0.038076   3rd Qu.: 0.05068   3rd Qu.: 0.031248   3rd Qu.: 0.035644  
##  Max.   : 0.110727   Max.   : 0.05068   Max.   : 0.170555   Max.   : 0.132044  
##        tc                 ldl                 hdl           
##  Min.   :-0.126781   Min.   :-0.115613   Min.   :-0.102307  
##  1st Qu.:-0.034248   1st Qu.:-0.030358   1st Qu.:-0.035117  
##  Median :-0.004321   Median :-0.003819   Median :-0.006584  
##  Mean   : 0.000000   Mean   : 0.000000   Mean   : 0.000000  
##  3rd Qu.: 0.028358   3rd Qu.: 0.029844   3rd Qu.: 0.029312  
##  Max.   : 0.153914   Max.   : 0.198788   Max.   : 0.181179  
##       tch                 ltg                 glu                target     
##  Min.   :-0.076395   Min.   :-0.126097   Min.   :-0.137767   Min.   : 25.0  
##  1st Qu.:-0.039493   1st Qu.:-0.033249   1st Qu.:-0.033179   1st Qu.: 87.0  
##  Median :-0.002592   Median :-0.001948   Median :-0.001078   Median :140.5  
##  Mean   : 0.000000   Mean   : 0.000000   Mean   : 0.000000   Mean   :152.1  
##  3rd Qu.: 0.034309   3rd Qu.: 0.032433   3rd Qu.: 0.027917   3rd Qu.:211.5  
##  Max.   : 0.185234   Max.   : 0.133599   Max.   : 0.135612   Max.   :346.0
\end{verbatim}

Bộ dữ liệu gồm 10 biến giải thích và 1 biến phụ thuộc \texttt{target}.
Các biến giải thích đã được chuẩn hóa nên có giá trị trung bình xấp xỉ
0. Biến \texttt{target} dao động từ 25 đến 346, với trung bình khoảng
152.1, cho thấy mức độ tiến triển bệnh tiểu đường có sự khác biệt đáng
kể giữa các bệnh nhân.

Trước khi lựa chọn biến, báo cáo tiến hành khám phá dữ liệu để quan sát
mối quan hệ giữa biến phụ thuộc \texttt{target} và các biến giải thích

\begin{Shaded}
\begin{Highlighting}[]
\FunctionTok{library}\NormalTok{(ggplot2)}
\FunctionTok{library}\NormalTok{(tidyr)}

\NormalTok{df }\SpecialCharTok{\%\textgreater{}\%}
  \FunctionTok{pivot\_longer}\NormalTok{(}\AttributeTok{cols =} \SpecialCharTok{{-}}\NormalTok{target, }\AttributeTok{names\_to =} \StringTok{"variable"}\NormalTok{, }\AttributeTok{values\_to =} \StringTok{"value"}\NormalTok{) }\SpecialCharTok{\%\textgreater{}\%}
  \FunctionTok{ggplot}\NormalTok{(}\FunctionTok{aes}\NormalTok{(}\AttributeTok{x =}\NormalTok{ value, }\AttributeTok{y =}\NormalTok{ target)) }\SpecialCharTok{+}
  \FunctionTok{geom\_point}\NormalTok{(}\AttributeTok{alpha =} \FloatTok{0.5}\NormalTok{) }\SpecialCharTok{+}
  \FunctionTok{geom\_smooth}\NormalTok{(}\AttributeTok{method =} \StringTok{"lm"}\NormalTok{, }\AttributeTok{se =} \ConstantTok{FALSE}\NormalTok{) }\SpecialCharTok{+}
  \FunctionTok{facet\_wrap}\NormalTok{(}\SpecialCharTok{\textasciitilde{}}\NormalTok{ variable, }\AttributeTok{scales =} \StringTok{"free\_x"}\NormalTok{) }\SpecialCharTok{+}
  \FunctionTok{labs}\NormalTok{(}
    \AttributeTok{title =} \StringTok{"Quan hệ giữa các biến của mô hình và target"}\NormalTok{,}
    \AttributeTok{x =} \StringTok{"Biến"}\NormalTok{,}
    \AttributeTok{y =} \StringTok{"target"}
\NormalTok{  )}
\end{Highlighting}
\end{Shaded}

\begin{verbatim}
## `geom_smooth()` using formula = 'y ~ x'
\end{verbatim}

\pandocbounded{\includegraphics[keepaspectratio]{DuLieuVaThietLap_files/DuLieuVaThietLap_files/figure-latex/unnamed-chunk-3-1.pdf}}

Từ các biểu đồ phân tán ở trên có thể thấy các biến \texttt{bmi},
\texttt{ltg}, \texttt{map} và \texttt{glu} có xu hướng tăng cùng với
\texttt{target}, trong khi \texttt{hdl} có xu hướng giảm khi
\texttt{target} tăng. Trong số đó, mối quan hệ của \texttt{bmi} và
\texttt{ltg} với \texttt{target} thể hiện khá rõ qua dạng phân bố của
các điểm dữ liệu.

Các biến \texttt{age}, \texttt{ldl}, \texttt{tc} và \texttt{tch} cũng
cho thấy xu hướng đồng biến với \texttt{target}, tuy nhiên mức độ không
rõ bằng các biến trên. Riêng biến \texttt{sex} gần như không thể hiện
mối quan hệ tuyến tính đáng kể với \texttt{target} khi xét riêng lẻ.

Nhìn chung, các scatter plots giúp nhận diện sơ bộ những biến có khả
năng ảnh hưởng đến \texttt{target} trước khi tiến hành xây dựng và lựa
chọn mô hình hồi quy.

\begin{Shaded}
\begin{Highlighting}[]
\NormalTok{cor\_matrix }\OtherTok{\textless{}{-}} \FunctionTok{cor}\NormalTok{(df)}

\FunctionTok{corrplot}\NormalTok{(}
\NormalTok{  cor\_matrix,}
  \AttributeTok{method =} \StringTok{"color"}\NormalTok{,}
  \AttributeTok{type =} \StringTok{"upper"}\NormalTok{,}
  \AttributeTok{tl.cex =} \FloatTok{0.8}\NormalTok{,}
  \AttributeTok{addCoef.col =} \StringTok{"black"}\NormalTok{,}
  \AttributeTok{number.cex =} \FloatTok{0.6}
\NormalTok{)}
\end{Highlighting}
\end{Shaded}

\pandocbounded{\includegraphics[keepaspectratio]{DuLieuVaThietLap_files/DuLieuVaThietLap_files/figure-latex/unnamed-chunk-4-1.pdf}}

Ma trận tương quan cho thấy biến target có tương quan dương tương đối rõ
với các biến bmi, map, tc, ldl, tch, ltg, glu. Trong đó, ltg và bmi là
hai biến có tương quan cao nhất với target. Ngược lại, hdl có tương quan
âm với target, còn sex gần như không có tương quan tuyến tính đơn biến
với target.

Một số cặp biến độc lập có tương quan cao, đặc biệt là tc-ldl,hdl-tch,
ldl-tch, tch-ltg cho thấy có thể tồn tại hiện tượng đa cộng tuyến. Do
đó, việc lựa chọn biến là cần thiết để xây dựng mô hình gọn hơn, tránh
sử dụng nhiều biến mang thông tin trùng lặp.

\subsection{Xây dựng mô hình hồi quy tuyến tính đầy đủ
biến}\label{xuxe2y-dux1ef1ng-muxf4-huxecnh-hux1ed3i-quy-tuyux1ebfn-tuxednh-ux111ux1ea7y-ux111ux1ee7-biux1ebfn}

Để đánh giá các phương pháp lựa chọn biến và lựa chọn mô hình, bộ dữ
liệu được chia thành hai phần: tập huấn luyện (train) và tập kiểm tra
(test) theo tỷ lệ 80:20. Tập train được sử dụng để xây dựng mô hình và
thực hiện lựa chọn biến, trong khi tập test được dùng để đánh giá khả
năng dự báo của các mô hình sau khi lựa chọn.

\begin{Shaded}
\begin{Highlighting}[]
\FunctionTok{set.seed}\NormalTok{(}\DecValTok{347}\NormalTok{)}

\CommentTok{\# Chia train:test theo tỉ lệ 80:20}
\NormalTok{train\_index }\OtherTok{\textless{}{-}} \FunctionTok{sample}\NormalTok{(}\DecValTok{1}\SpecialCharTok{:}\FunctionTok{nrow}\NormalTok{(df), }\AttributeTok{size =} \FloatTok{0.8} \SpecialCharTok{*} \FunctionTok{nrow}\NormalTok{(df))}

\NormalTok{train\_data }\OtherTok{\textless{}{-}}\NormalTok{ df[train\_index, ]}
\NormalTok{test\_data  }\OtherTok{\textless{}{-}}\NormalTok{ df[}\SpecialCharTok{{-}}\NormalTok{train\_index, ]}

\FunctionTok{dim}\NormalTok{(train\_data)}
\end{Highlighting}
\end{Shaded}

\begin{verbatim}
## [1] 353  11
\end{verbatim}

\begin{Shaded}
\begin{Highlighting}[]
\FunctionTok{dim}\NormalTok{(test\_data)}
\end{Highlighting}
\end{Shaded}

\begin{verbatim}
## [1] 89 11
\end{verbatim}

\begin{Shaded}
\begin{Highlighting}[]
\NormalTok{full\_model }\OtherTok{\textless{}{-}} \FunctionTok{lm}\NormalTok{(target }\SpecialCharTok{\textasciitilde{}}\NormalTok{ ., }\AttributeTok{data =}\NormalTok{ train\_data)}

\FunctionTok{summary}\NormalTok{(full\_model)}
\end{Highlighting}
\end{Shaded}

\begin{verbatim}
## 
## Call:
## lm(formula = target ~ ., data = train_data)
## 
## Residuals:
##      Min       1Q   Median       3Q      Max 
## -153.361  -37.612   -1.205   37.236  148.804 
## 
## Coefficients:
##             Estimate Std. Error t value Pr(>|t|)    
## (Intercept)  153.437      2.923  52.495  < 2e-16 ***
## age           12.981     68.640   0.189 0.850109    
## sex         -255.072     68.749  -3.710 0.000242 ***
## bmi          478.669     75.547   6.336 7.43e-10 ***
## map          355.903     73.341   4.853 1.85e-06 ***
## tc          -749.672    495.989  -1.511 0.131592    
## ldl          497.297    398.898   1.247 0.213369    
## hdl           55.853    250.479   0.223 0.823680    
## tch           52.106    188.380   0.277 0.782254    
## ltg          766.194    196.771   3.894 0.000119 ***
## glu          109.987     73.685   1.493 0.136449    
## ---
## Signif. codes:  0 '***' 0.001 '**' 0.01 '*' 0.05 '.' 0.1 ' ' 1
## 
## Residual standard error: 54.82 on 342 degrees of freedom
## Multiple R-squared:  0.5075, Adjusted R-squared:  0.4931 
## F-statistic: 35.24 on 10 and 342 DF,  p-value: < 2.2e-16
\end{verbatim}

Mô hình hồi quy đầy đủ có ý nghĩa thống kê tổng thể với p-value của kiểm
định F nhỏ hơn 2.2e-16. Giá trị \(R^2 = 0.5075\) cho thấy mô hình giải
thích được khoảng 50.75\% biến thiên của \texttt{target}, trong khi
\(R^2_{adj} = 0.4931\) cho thấy sau khi điều chỉnh theo số lượng biến,
tỷ lệ giải thích còn khoảng 49.31\%.

Xét từng biến giải thích, \texttt{sex}, \texttt{bmi}, \texttt{map} và
\texttt{ltg} có ý nghĩa thống kê ở mức 5\%. Trong đó, \texttt{bmi},
\texttt{map} và \texttt{ltg} có ảnh hưởng cùng chiều đến
\texttt{target}, còn \texttt{sex} có ảnh hưởng ngược chiều. Các biến còn
lại có p-value lớn hơn 0.05 nên chưa có đủ bằng chứng để khẳng định ảnh
hưởng riêng lẻ của chúng trong mô hình khi đã xét đồng thời các biến
khác.

Mặc dù mô hình đầy đủ sử dụng toàn bộ 10 biến giải thích, chỉ một số
biến thực sự có ý nghĩa thống kê. Điều này cho thấy mô hình có thể chứa
những biến không cần thiết. Vì vậy, các phương pháp lựa chọn biến sẽ
được áp dụng ở các phần tiếp theo nhằm tìm ra mô hình gọn hơn nhưng vẫn
đảm bảo khả năng giải thích và dự báo.

\end{document}
